% ATTEMPT_STATUS: unresolved
% RESEARCH_GENERATED: 2026-10-01; GPT-6 Astra Ultra; individual mathematical attempt
% TOP_PROBLEM: {"id": 160, "title": "Large Subsets of Approximate Groups", "category_id": 4, "category": {"id": 4, "name": "algebra", "display_name": "Algebra", "description": "Group theory, ring theory, field theory, and algebraic structures.", "slug": "algebra", "order_index": 4, "created_at": "2026-07-31T15:26:25.671Z"}, "status": "open"}
% FIRST_POSTED: 2026-10-01
% Imported from: unsolvedmath_additional_500.tex; UnsolvedMath ID 160
% !TeX program = lualatex
% Compile twice: lualatex 160.tex
% Self-contained: no external bibliography, images, or \input files.
% Numeric section labels and visible numbers are original UnsolvedMath IDs.
\documentclass[11pt,a4paper]{article}
\usepackage[margin=24mm,headheight=14pt]{geometry}
\usepackage{fontspec}
\setmainfont{Latin Modern Roman}
\setsansfont{Latin Modern Sans}
\setmonofont{Latin Modern Mono}
\usepackage{amsmath,amssymb,mathtools}
\usepackage{unicode-math}
\setmathfont{Latin Modern Math}
\usepackage{microtype}
\usepackage{enumitem,xurl,needspace}
\usepackage[dvipsnames]{xcolor}
\usepackage{hyperref}
\hypersetup{unicode=true,colorlinks=true,linkcolor=MidnightBlue,urlcolor=MidnightBlue,
 pdftitle={UnsolvedMath Problem 160},pdfauthor={Source-annotated compilation},bookmarksnumbered=true}
\usepackage{bookmark,tocloft,titlesec,fancyhdr}
\setlength{\cftsecnumwidth}{4.2em}
\cftsetpnumwidth{2.5em}
\cftsetrmarg{4em}
\titleformat{\section}{\Large\bfseries\raggedright}{\thesection}{0.7em}{}
\pagestyle{fancy}\fancyhf{}
\fancyhead[L]{\small\sffamily UnsolvedMath research attempt}
\fancyhead[R]{\small Original catalogue IDs}
\fancyfoot[C]{\thepage}
\setlength{\parindent}{0pt}
\setlength{\parskip}{5pt plus 1pt}
\setlength{\emergencystretch}{3em}
\setcounter{tocdepth}{1}\setcounter{secnumdepth}{1}
\setlist[itemize]{leftmargin=3em,itemsep=4pt,topsep=4pt,parsep=0pt}
\allowdisplaybreaks
\newcommand{\N}{\mathbb{N}}
\newcommand{\Z}{\mathbb{Z}}
\newcommand{\Q}{\mathbb{Q}}
\newcommand{\R}{\mathbb{R}}
\newcommand{\C}{\mathbb{C}}
\newcommand{\F}{\mathbb{F}}
\newcommand{\A}{\mathbb{A}}
\newcommand{\PP}{\mathbb{P}}
\newcommand{\E}{\mathbb{E}}
\newcommand{\eps}{\varepsilon}
\newcommand{\dd}{\,\mathrm d}
\newcommand{\sref}[2]{\hyperlink{source:#1:#2}{[S#2]}}
\newif\ifshowcatalogue
\showcataloguetrue
\newcommand{\cataloguescope}[1]{\ifshowcatalogue
 \paragraph{Statement recorded in the supplied source.}
 #1\fi}
\begin{document}
% UnsolvedMath ID 160; source catalogue code GREEN-072

\setcounter{section}{159}

\section{Polynomial-Size Subsets with Controlled Products}\label{160}

{\small\sffamily UnsolvedMath ID 160 · Algebra}

\subsection{Definitions and mathematical statement}

\paragraph{Shared notation.}Unless a section says otherwise, $\N=\{1,2,\ldots\}$,
$\Z$ is the ring of integers, $\Q,\R,\C$ are the rational, real, and complex
fields, $[N]=\{1,\ldots,\lfloor N\rfloor\}$, and $\log$ is the natural
logarithm. For real functions of a parameter tending to infinity,
$f\ll g$ and $f=O(g)$ mean $|f|\leq Cg$ eventually for some fixed $C$;
$f\gg g$ reverses this comparison; $f=o(g)$ means $f/g\to0$;
$f\sim g$ means $f/g\to1$; and $f\asymp g$ means both order bounds.
A subscript on $O$ or $\ll$ specifies allowed constant dependence.
The expression $n^{o(1)}$ in an upper bound means growth smaller than
$n^\epsilon$ for each fixed $\epsilon>0$. Local definitions govern any
exception, finite-field convention, or use of the same symbol for another
object. An asymptotic claim is not automatically an effective algorithm.



A finite $K$-approximate group is a finite subset $A$ of a group with $1\in A$, $A=A^{-1}$, and $A^2\subseteq XA$ for some $X$ with $|X|\leq K$. For $j\geq1$, $A^j$ consists of products of $j$ elements, with repetitions allowed. Seek absolute constants $c,C>0$ and $S\subseteq A$ with $|S|\geq cK^{-C}|A|$ and $S^8\subseteq A^4$. \sref{160}{1} \sref{160}{2}

\paragraph{Statement recorded in the supplied source.}
Suppose $A$ is a $K$-approximate group (not necessarily abelian). Is there $S \subset A$ with $|S| \gg K^{-O(1)}|A|$ and $S^8 \subset A^4$? \sref{160}{1}

\subsection{Short English statement}

Can every finite approximate group contain a polynomially large subset whose eightfold products stay inside the original fourfold product set?

\subsection{Research attempt}

\paragraph{Provenance and scope.}
This individual attempt was generated by GPT-6 Astra Ultra on 1 October 2026. The method was to derive explicit partial results and test the first proposed extension adversarially. The original statement and sources below are retained. No claim of mathematical novelty or complete solution is made.

\paragraph{Formal target and context.}
A finite $K$-approximate group is symmetric, contains the identity, and satisfies $A^2\subset XA$ with $|X|\leq K$. The desired subset need not be a subgroup or symmetric; it must have polynomial relative size in $K$ and satisfy the strong containment $S^8\subset A^4$. The inspected Breuillard--Green--Tao survey supplies the surrounding approximate-group context. We prove the target for a concrete full-rank box family with its optimal covering parameter, and test two shortcuts that fail even in the integers.

\paragraph{The exact $K=1$ case.}
If $A^2\subset xA$ for one translate, then $|A^2|\leq|A|$. Since the identity belongs to $A$, also $A\subset A^2$, so equality holds and $A^2=A$. Symmetry and the identity then make $A$ a subgroup. Taking $S=A$ gives $S^8=A\subset A^4$. Thus the first parameter value is settled exactly, with no loss in size. The argument uses finiteness in the cardinality step.

\paragraph{Integer boxes have a polynomial-size inner set.}
In the additive group $\Z^d$, take
\[
 A=\prod_{i=1}^d[-L_i,L_i]\cap\Z^d,\qquad L_i\geq1.
\]
The doubled box is covered by the $2^d$ translates $x+A$ with $x_i\in\{-L_i,L_i\}$. Conversely each of the $2^d$ corner points of $2A$, whose coordinates are $\pm2L_i$, requires a different translate of $A$: two different corners differ by $4L_i$ in some coordinate, whereas a translate of $A$ has coordinate width $2L_i$. Thus its least covering parameter is exactly $K=2^d$.

Let
\[
 S=\prod_{i=1}^d[-\lfloor L_i/2\rfloor,\lfloor L_i/2\rfloor]\cap\Z^d.
\]
Eight summands from $S$ have coordinate absolute value at most $8\lfloor L_i/2\rfloor\leq4L_i$, so $8S\subset4A$. Moreover
\[
 \frac{|S|}{|A|}=
 \prod_i\frac{2\lfloor L_i/2\rfloor+1}{2L_i+1}
 \geq3^{-d}=K^{-\log_2 3}.
\]
The one-dimensional ratio is at least one third, with equality at $L_i=1$; checking even and odd $L_i$ proves the bound for all positive integer side lengths. Therefore the exact polynomial-size conclusion holds for every such box, including thin boxes with all side lengths one.

The same containment survives a homomorphism into an abelian group. To retain the displayed cardinality ratio, however, injectivity on the original box is needed; collisions can otherwise change numerator and denominator differently. We do not silently extend the rank calculation to arbitrary nonproper progressions or noncommuting product sets.

\paragraph{Why searching only for a subgroup would miss valid examples.}
For $A=[-L,L]\subset\Z$, the covering parameter is two, uniformly in $L$. The only finite subgroup of $\Z$ is $\{0\}$. Any method insisting that $S$ itself be a subgroup therefore gives size one and cannot meet a fixed positive fraction of $|A|$ as $L$ grows. Nevertheless the half-interval construction above does meet the desired containment and has size comparable to $|A|$. The original question deliberately allows this kind of inner region; replacing it by a subgroup problem is strictly stronger and already false for this family.

\paragraph{Small growth does not by itself prove containment.}
From $A^2\subset XA$, associativity gives $A^j\subset X^{j-1}A$ by induction and therefore $|A^j|\leq K^{j-1}|A|$. These are cardinality bounds. They do not imply $A^8\subset A^4$, and selecting $S=A$ is invalid in general: for the integer interval, $8A=[-8L,8L]$ while $4A=[-4L,4L]$. A proof must locate elements whose repeated products stay in the smaller set, not merely compare the sizes of two larger powers.

In the box example, coordinates give an explicit notion of shrinking. In an arbitrary nonabelian approximate group, multiplication order and commutators prevent copying that coordinate argument without additional structure. The missing information is a sufficiently large controlled inner set, not just a polynomial bound for growth of a fixed number of powers.

\paragraph{Verification and remaining gap.}
All box powers are ordinary repeated Minkowski sums of full integer intervals, so their endpoint descriptions are exact. The covering parameter was proved minimal rather than assigned artificially, and the size loss is therefore polynomial in the actual parameter. The subgroup counterexample respects every approximate-group hypothesis. The attempt proves the full target for $K=1$ and for integer boxes, while exposing two invalid general reductions. The uniform nonabelian statement remains unresolved here.

\subsection{Sources}

{\small

\begin{itemize}

\item[\hypertarget{source:160:1}{S1}] ULAM.AI, \emph{UnsolvedMath}, supplied \texttt{problems.json}, record 160 (GREEN-072), statement and background. The embedded literature-review date is 2026-08-17; that is the archive’s review date, not a fresh certification. Dataset landing page: \url{https://huggingface.co/datasets/ulamai/UnsolvedMath}.

\item[\hypertarget{source:160:2}{S2}] Emmanuel Breuillard, Ben Green and Terence Tao, Small doubling in groups, arXiv:1301.7718; Erdos Centennial, Bolyai Society Mathematical Studies 25 (2013), 129-151. \url{https://arxiv.org/abs/1301.7718}. Bibliographic information imported from the supplied record.

\item[\hypertarget{source:160:3}{S3}] Ben Green, 100 Open Problems, Problem 29, most recent update December 2025. \url{https://people.maths.ox.ac.uk/greenbj/papers/open-problems.pdf}. The author-hosted document was inspected on 27 September 2026; the archive’s local code is not a problem-number locator in this paper.

\item[E1] E. Breuillard, B. Green and T. Tao, \emph{Small doubling in groups}. \url{https://arxiv.org/abs/1301.7718}. Inspected 1 October 2026.

\end{itemize}

}

\label{end:160}
\end{document}
